% Generated draft: requires mathematical and source review.
\documentclass[11pt]{article}
\usepackage[margin=25mm]{geometry}
\usepackage{fontspec}
\setmainfont{Latin Modern Roman}
\usepackage{amsmath,amssymb,mathtools,unicode-math}
\setmathfont{Latin Modern Math}
\usepackage{enumitem,needspace,xurl,hyperref}
\setlength{\parindent}{0pt}
\setlength{\parskip}{5pt}
\setlength{\emergencystretch}{3em}
\newcommand{\N}{\mathbb{N}}
\newcommand{\Z}{\mathbb{Z}}
\newcommand{\Q}{\mathbb{Q}}
\newcommand{\R}{\mathbb{R}}
\newcommand{\C}{\mathbb{C}}
\newcommand{\F}{\mathbb{F}}
\newcommand{\A}{\mathbb{A}}
\newcommand{\PP}{\mathbb{P}}
\newcommand{\E}{\mathbb{E}}
\newcommand{\eps}{\varepsilon}
\newcommand{\dd}{\,\mathrm{d}}
\begin{document}
\section*{UnsolvedMath ID: 8800022}
\textbf{Source problem code:} catalogue entry
\subsection{Definitions and mathematical statement}
An integer factoring task takes a composite integer \(N\) and seeks a divisor \(d\) with \(1<d<N\). In a partial-information attack, some bits of an associated secret integer are supplied, and additional unknown bits may be guessed before a factoring procedure is attempted. If that integer is a factor \(p\mid N\), a wrong guess is an assignment disagreeing with the true binary digits of \(p\); the source does not say that factor bits are the only intended information.

May's task is to extract useful information from such incorrect guesses, rather than simply discard the corresponding attempts. A proposed solution would have to identify what data an unsuccessful attempt produces and how those data yield valid constraints, better guesses, or a justified reduction of the remaining work needed to factor \(N\).

The excerpt supplies no attack algorithm, bit positions or counts, available transcripts, probability distribution, or required improvement. A failed factoring attempt is not automatically a certificate that the guessed bits were wrong: it may fail even on a correct guess unless the underlying procedure has a stated success guarantee. Thus neither a correctness oracle for guesses nor an information-theoretic learning model is silently assumed. The task remains an open-ended algorithmic construction problem with the intended attack setting missing.
\subsection{Short English statement}
Can unsuccessful guessed-bit attempts in partial-information factoring attacks provide information that improves later attempts? The attack, observable failure data, and the meaning of useful improvement need specification.
\subsection{Sources}
\begin{itemize}
\item Local UnsolvedMath registry, permanent ID 8800022, source code catalogue entry; attributed to Alexander May; no proposal year supplied; update 2026-08-06.
\item IPAM, \emph{Workshop on Number Theory and Cryptography --- Open Problems} (2006), attachment 6222, PDF page 13, printed slide 9: \url{https://www.ipam.ucla.edu/programs/workshops/workshop-i-number-theory-and-cryptography-open-problems/?tab=schedule}. Registry access date 2026-07-29; PDF and legacy-PPT visual audit. Provenance reports 19 attachments, 17 PDFs and two legacy PPT decks, with repetitions and solved definitions excluded. Present primary-literature status and redistribution terms require review.
\item The supplied summary describes an unresolved open-ended question in May's slides and reports no verified specific solution or substantial progress. No additional exact reference is supplied; this is not an independent literature assessment.
\end{itemize}
\end{document}
